% ===================================================================
%  TABELO N-RO 15 — La bazaro. — La manĝeblaĵoj.
%  Traduction 2026 d'apres texto/io, controlee sur texto/fr et
%  texto/es. Consigne : voir texto/eo/10-tabelo-01.tex.
% ===================================================================

%%K t15-apar-1 apar
\VUsaut{4.0mm}
\VUtitre{15.0pt}{60}{TABELO N\textsuperscript{o} 15}
\VUsaut{3.0mm}

%%K t15-tit-1 sub
\VUpk{11.4pt}{40}{La bazaro. --- La manĝeblaĵoj.}
\VUsaut{3.0mm}

%%K t15-01-1 p
1. --- La plej multaj komercistoj kiuj liveras al ni la varojn necesajn por nia
nutrado sin kunigas sub la \VUgras{halo} \textsuperscript{(1)} de la
\VUgras{kovrita bazaro} aŭ en la najbaraj stratoj.

%%K t15-02-1 p
2. --- La \VUgras{porkoviandisto} \textsuperscript{(2)}, kiu vendas
\VUgras{ŝinkon} \textsuperscript{(3)}, \VUgras{sangokolbason}
\textsuperscript{(4)}, \VUgras{kolbaseton} \textsuperscript{(5)},
\VUgras{tripkolbaseton} \textsuperscript{(6)} kaj \VUgras{lardon}
\textsuperscript{(7)}, staras apud la \VUgras{buter- kaj fromaĝvendistino}
\textsuperscript{(8)}, kies elmetejo estas provizita per \VUgras{holandaj
fromaĝoj} \textsuperscript{(9)} kun ruĝa krusto, per \VUgras{kamember-fromaĝoj}
\textsuperscript{(10)}, per \VUgras{gruyère-diskoj} \textsuperscript{(11)} kaj
per freŝaj \VUgras{buterblokoj} \textsuperscript{(12)}.

%%K t15-03-1 p
3. --- Antaŭ ŝi, \VUgras{legomvendistino} \textsuperscript{(13)} proponas al
siaj klientinoj belajn \VUgras{tomatojn} \textsuperscript{(14)},
\VUgras{florbrasikojn} \textsuperscript{(15)}, \VUgras{salaton}
\textsuperscript{(16)}, \VUgras{kapbrasikojn} \textsuperscript{(17)},
\VUgras{kukurbetojn} \textsuperscript{(19)}, \VUgras{melonojn}
\textsuperscript{(18)} kaj eĉ \VUgras{fungojn} \textsuperscript{(20)}. Sed
\VUgras{bierkelnero}, kiu haltigis sian \VUgras{liverveturilon}
\textsuperscript{(21)} ŝarĝitan per \VUgras{bierbareloj}
\textsuperscript{(22)}, ĝuste antaŭ ŝia elmetejo, malhelpas ŝiajn klientojn
alproksimiĝi. Ĝardenistino kunportas al ŝi korbegon da \VUgras{artiŝokoj}
\textsuperscript{(23)}.

%%K t15-tit-2 sub
\VUsaut{4.0mm}
\VUpk{10.2pt}{40}{Vendado kaj aĉetado.}
\VUsaut{3.0mm}

%%K t15-04-1 p
4. --- En la bazaro, la vivo estas granda, ĉar alvenis la horo de la
\VUgras{aŭkcia vendo} \textsuperscript{(24)}. La \VUgras{domomastrinoj}
\textsuperscript{(26)}, kiuj venas tien por siaj aĉetoj, haltas pasante antaŭ la
butiko de la \VUgras{tripvendistino} \textsuperscript{(25)}, por aĉeti
\VUgras{stomaktripojn} aŭ \VUgras{bovidkapon}.

%%K t15-05-1 p
5. --- Grasa \VUgras{kuiristo} \textsuperscript{(27)} marĉandas tre belan
\VUgras{tinuson} ĉe la \VUgras{fiŝvendistino} \textsuperscript{(28)} kiu laŭ sia
arbitro pli altigas siajn (vend)prezojn. Ŝi ricevis grandan kvanton da
\VUgras{angiloj, soleoj, trutoj} kaj \VUgras{karpoj}. Ŝia \VUgras{moruo} kaj
\VUgras{mituloj} estas tre freŝaj.

%%K t15-tit-3 sub
\VUsaut{4.0mm}
\VUpk{10.2pt}{40}{Malkaraj kaj karaj varoj.}
\VUsaut{3.0mm}

%%K t15-06-1 p
6. --- La \VUgras{kortbirdvendistino} \textsuperscript{(29)}, instalita apud ŝi,
estas tre konscienca. Kiam ŝi vendas \VUgras{meleagron, anseron, anason} aŭ
simplan \VUgras{kokidon}, ŝi postulas nur la justan prezon.

%%K t15-07-1 p
7. --- Ĉe la angulo de la placo, sur la unua etaĝo, de mallonga tempo, estas
instalita \VUgras{tolvendistino} \textsuperscript{(30)} ĉe kiu la aĉetantinoj
estas ĉiam certaj trovi tre belan sortimenton da \VUgras{tablotukoj,
buŝtukoj, antaŭtukoj} kaj \VUgras{viŝiloj}. Sur la sama etaĝo,
\VUgras{tranĉilisto} \textsuperscript{(31)} instalis sian laborejon. Li akrigas
la \VUgras{tranĉilojn} de la vendistinoj en la halo kaj fabrikas por la
ĝardenistoj \VUgras{serpojn} el la plej malmola \VUgras{ŝtalo}.

%%K t15-tit-4 sub
\VUsaut{4.0mm}
\VUpk{10.2pt}{40}{La spicejo.}
\VUsaut{3.0mm}

%%K t15-08-1 p
8. --- Sub lia laborejo, de mallonga tempo, estas malfermita granda magazeno da
nutraĵoj. Jam ĝi ne plu estas la simpla kaj negrava \VUgras{spicvendejo}
\textsuperscript{(32)} iama, kiu nur vendis la komunuzajn varojn, \VUgras{teon}
\textsuperscript{(33)}, \VUgras{ĉokoladon} \textsuperscript{(34)},
\VUgras{blokan sukeron} \textsuperscript{(37)} kaj \VUgras{sapon}
\textsuperscript{(38)}. Tie oni vendas iun ajn varon, \VUgras{krakantajn
kukojn}, \VUgras{konservaĵojn} \textsuperscript{(35)}, \VUgras{likvorojn}
\textsuperscript{(36)}, \VUgras{rumon, konjakon, kirŝon, anizlikvoron} kaj
\VUgras{kuracaon}. Tie oni trovas ĉasaĵon: \VUgras{leporon}
\textsuperscript{(39)}, \VUgras{turdojn} \textsuperscript{(40)},
\VUgras{perdrikojn} \textsuperscript{(41)}, \VUgras{koturnojn}
\textsuperscript{(42)}, \VUgras{sovaĝan anason} \textsuperscript{(43)} aŭ
\VUgras{fazanon} \textsuperscript{(44)}, tiel facile kiel ekzotikajn fruktojn
tiajn kiaj \VUgras{bananoj} \textsuperscript{(46)} kaj \VUgras{ananasoj}.

%%K t15-09-1 p
9. --- Spicista \VUgras{metilernanto} \textsuperscript{(45)} tre komplezema,
estas preta por doni tion kion oni deziras: \VUgras{konfitaĵon}
\textsuperscript{(47)} riban aŭ cidonian, \VUgras{grason} \textsuperscript{(48)},
\VUgras{kukumetojn} \textsuperscript{(49)} aŭ salitajn \VUgras{sardinojn}
\textsuperscript{(50)}.

%%K t15-09-2 p
Tio estas tre oportuna por la \VUgras{klientinoj} \textsuperscript{(51)}, kiuj
trovas, apud la plej komunuzaj varoj, la fruktunuaĵojn plej maloftajn kaj la
fruktojn plej bongustajn: sukoplenajn \VUgras{pirojn} \textsuperscript{(52)},
\VUgras{prunojn} \textsuperscript{(53)}, \VUgras{vinberojn}
\textsuperscript{(54)}, \VUgras{persikojn} \textsuperscript{(55)},
\VUgras{migdalojn} \textsuperscript{(56)} kaj \VUgras{juglandojn}
\textsuperscript{(57)}.

%%K t15-tit-5 sub
\VUsaut{4.0mm}
\VUpk{10.2pt}{40}{La klientaro.}
\VUsaut{3.0mm}

%%K t15-10-1 p
10. --- La \VUgras{rizo} \textsuperscript{(58)} kaj la \VUgras{lentoj}
\textsuperscript{(59)} estas apud la kesto de la fumaĵitaj \VUgras{haringoj}
\textsuperscript{(60)}; la \VUgras{terpomo} \textsuperscript{(61)} kaj la
\VUgras{fazeolo} \textsuperscript{(63)} estas najbaraj al la \VUgras{pomoj}
\textsuperscript{(62)}, \VUgras{figoj} \textsuperscript{(64)}, \VUgras{sekaj
prunoj} \textsuperscript{(65)} kaj \VUgras{aveloj} \textsuperscript{(66)}. Oni
trovas en tiu magazeno freŝajn \VUgras{ovojn} \textsuperscript{(67)} kaj
\VUgras{olivojn} \textsuperscript{(68)}; pro tio \VUgras{korboj}
\textsuperscript{(70)} kaj \VUgras{provizretoj} \textsuperscript{(73)} estas
rapidege plenigitaj.

%%K t15-11-1 p
11. --- La \VUgras{kafo} \textsuperscript{(72)} de tiu bonega firmo akiris,
inter la \VUgras{servistinoj} \textsuperscript{(69)} kaj la
\VUgras{kuiristinoj}, justan reputacion, ĉar la maljuna \VUgras{spicisto}
\textsuperscript{(71)} mem plej zorge rostas ĝin.

%%K t15-tit-6 sub
\VUsaut{4.0mm}
\VUpk{10.2pt}{40}{La spektaĵoj.}
\VUsaut{3.0mm}

%%K t15-12-1 p
12. --- Ne ĉiuj iras al la bazaro por provizi sin. En la pentrinda cirkulado
sur la strateto kiu kondukas al la halo, oni vidas la plej diversajn personojn.
Apud afabla \VUgras{buĉservisto} \textsuperscript{(75)}, kiu puŝas antaŭ si
\VUgras{puŝĉareton} \textsuperscript{(74)}, jen du forpermesitaj soldatoj tute
fieraj vidigi sian brilantan uniformon: la \VUgras{husaro}
\textsuperscript{(76)} kun blua \VUgras{dolmano} kaj \VUgras{plumhava ĉako}, kaj
la \VUgras{zuava kaporalo}, kun pantalono pufa kaj kapkovrita per \VUgras{fezo}.

%%K t15-13-1 p
13. --- La \VUgras{panisto} \textsuperscript{(80)}, kiu staras ĉe la sojlo de la
\VUgras{panejo} \textsuperscript{(78)} ĵus knedis la paston de sia \VUgras{pano}
\textsuperscript{(79)} kaj reprenis el la forno la \VUgras{kornopanojn}
\textsuperscript{(81)}, la \VUgras{bulkojn} \textsuperscript{(82)} kaj la
\VUgras{diskopanojn} \textsuperscript{(83)} kiuj estas en la vitrino.

%%K t15-14-1 p
14. --- Super la panejo loĝas lerta \VUgras{tolaĵistino} \textsuperscript{(84)}
kiu dungas plurajn \VUgras{brodistinojn} \textsuperscript{(85)}. Apud ŝi loĝas
\VUgras{lav-gladistino} \textsuperscript{(86)}, kiu amelas la \VUgras{tolaĵaron}
\textsuperscript{(88)} post lavi kaj sekigi ĝin, kaj kiu gladas ĝin per
\VUgras{gladilo} \textsuperscript{(87)}.

%%K t15-tit-7 sub
\VUsaut{4.0mm}
\VUpk{10.2pt}{40}{Viando, fiŝoj kaj legomoj.}
\VUsaut{3.0mm}

%%K t15-15-1 p
15. --- En la \VUgras{viandejo} \textsuperscript{(89)}, situanta sur la
teretaĝo, oni vendas ĉiaspecajn \VUgras{viandojn, ŝafviandon}
\textsuperscript{(90)}, \VUgras{bovviandon} \textsuperscript{(94)} aŭ
\VUgras{bovidviandon} \textsuperscript{(92)} kiel \VUgras{ŝafkrurojn,
bifstekojn, rostaĵojn} kaj \VUgras{kotletojn}. La \VUgras{buĉservisto}
\textsuperscript{(93)} dissekcas tiujn ĉiujn viandojn kaj aparte metas la
\VUgras{medolostojn} \textsuperscript{(96)}. Ĉe la pordo de la viandejo estas
\VUgras{ostrovendistino} \textsuperscript{(97)} kiu vendas \VUgras{ostrojn}
\textsuperscript{(98)}. Ŝi malfermas ilin se oni deziras tion.

%%K t15-16-1 p
16. --- Tiuj ĉiuj komercistoj pagas patenton aŭ lokimposton; pro tio la
\VUgras{policisto} \textsuperscript{(100)} senkompate persekutas la kompatindajn
\VUgras{legomkolportistinojn} \textsuperscript{(99)} kiuj konkurencas ilin
disportante per siaj veturiletoj \VUgras{karotojn, rafanetojn, napojn,
poreojn} aŭ \VUgras{asparagajn fasketojn, freŝajn pizojn, spinacojn, okzalon,
kreson} kaj \VUgras{petroselon}.

%%K t15-tit-8 sub
\VUsaut{4.0mm}
\VUpk{10.2pt}{40}{Atentu la policiston!}
\VUsaut{3.0mm}

%%K t15-17-1 p
17. --- La knabo kiu vendas \VUgras{ajlon} \textsuperscript{(101)} kaj
\VUgras{cepojn}, tre bone scias, ke ankaŭ li, ne rajtas vendi siajn varojn apud
la bazaro. Li forkuras, kurante kun risko renversi la korbon da \VUgras{kraboj},
\VUgras{kankroj} kaj \VUgras{salikokoj} de la \VUgras{vendistino}
\textsuperscript{(102)} de \VUgras{langustoj} kaj \VUgras{omaroj}.

%%K t15-18-1 p
18. --- Ĉiuj agitas sin kaj bruegas; sed tio ne malhelpas aŭdi klare la voĉon de
la kolportanta \VUgras{vitristo} \textsuperscript{(103)}, nek la krion de la
\VUgras{karbisto} \textsuperscript{(104)} kiu ĵus lasis dum kelkaj momentoj
sian ĉareton por liveri \VUgras{karban sakegon} \textsuperscript{(105)}.
\VUsaut{6.0mm}
\VUfilet{22mm}
